\section{Discussion}
\label{sec:discussion}

This section draws the experiments together and compares the techniques as answers to a
single question: how best to close the Sim2Real gap for runway keypoint detection under
adverse weather when training on nominal synthetic data.

\subsection{Comparing the Techniques}
\label{subsec:disc-compare}
The three intervention families differ sharply in effectiveness on this task.
\emph{Weather augmentation} (\Cref{sec:weather-aug}) is cheap and label-preserving but
delivers only limited, unstable gains: the scenario-matched augmentation is often not
the best, the useful ratio is condition-specific (e.g.\ water refraction $+10$ mAP on
the \texttt{ges} rain split at ratio $0.75$, but glare augmentation dropping $>10$ mAP
on the fog split from $0.45$ to $0.75$), and it does not recover the advantage of
source diversity. \emph{Self-training} (\Cref{sec:pseudo-labels}) is markedly stronger:
using unlabeled real footage it improves the adverse LARD splits by more than $10$ mAP
and the in-distribution landing sequences by up to $35$ mAP, with the confidence
threshold nearly neutral on LARD. \emph{Style transfer} (\Cref{sec:style-transfer})
is expected to complement these but is still ongoing. \Cref{tab:disc-summary}
consolidates the best configuration of each available technique against the
\texttt{combined} baseline.

\begin{table}[th]
    \caption{Best configuration of each technique vs.\ the \texttt{combined} baseline,
    real test set per weather scenario and overall (mAP@[.50:.95]). Style transfer is
    pending.}
    \label{tab:disc-summary}
    \centering
    \begin{NiceTabular}{lrrrrrr}
        \CodeBefore
        \rowcolors{2}{gray!10}{white}
        \Body
        \toprule
        \textbf{Technique} & \textbf{Nominal} & \textbf{Glare} & \textbf{Rain} & \textbf{Fog} & \textbf{Night} & \textbf{Overall} \\
        \midrule
        combined (baseline)            & 0.636 & 0.586 & 0.527 & 0.280 & 0.406 & \tbd \\
        + Weather augmentation (best)  & \tbd & \tbd & \tbd & \tbd & \tbd & \tbd \\
        + Self-training (best)         & \tbd & \tbd & \tbd & \tbd & \tbd & \tbd \\
        + Style transfer (Cosmos)      & \tbd & \tbd & \tbd & \tbd & \tbd & \tbd \\
        \bottomrule
    \end{NiceTabular}
\end{table}
\todo[inline]{Fill the best-configuration rows once the augmentation, self-training and
style-transfer runs are consolidated on a common (size+angle-filtered) real test set.}

\subsection{The Gap Is Real but Modest, and Partly a Quality Effect}
\label{subsec:disc-gap}
A recurring theme is that the raw real-versus-synthetic comparison is misleading. Only
after matching both runway size and vertical angle does the expected ordering appear
(each model best on its own source; real no longer the easiest set), and even then the
gap is modest. Much of the synthetic difficulty is an \emph{image-quality} effect ---
poor runway/background separability, confusion with parallel structures, and softer
edges on synthetic data --- most pronounced for distant runways. The practical
consequence is that cross-domain evaluations for this task must control for size and
orientation, or they will systematically flatter the real set.

\subsection{Dataset Diversity vs.\ Viewpoint Coverage}
\label{subsec:disc-diversity}
The \texttt{combined} source outperforms every individual source on real data despite
covering roughly one fifth as many locations, indicating that \emph{rendering-source
diversity outweighs runway-viewpoint coverage} in this setup. Consistent with the
LARD-V2 benchmark, no single source wins everywhere: satellite sources
(\texttt{ges}, \texttt{arcgis}, \texttt{bingmaps}) lead under glare and nominal, while
the 3D simulator \texttt{flsim} leads under rain, fog and night. The
augmentation study confirms the same lesson from the other direction: appearance
augmentation of a single source does not reach the diverse combined baseline.

\subsection{Limitations}
\label{subsec:disc-limitations}
Several limitations qualify these results. The weather-augmentation sweep uses a reduced
20-epoch budget and, even with a fixed training set, four runs do not fully average out
the sampling variance; larger effects are trustworthy, small ones less so. The
weather-stratified real splits --- especially fog and night --- are small, so
per-scenario numbers are noisier than the nominal figure. Self-training is exposed to
confirmation bias, as the rain pseudo-label composition shows (quantity $\neq$ quality),
and its largest gains are on in-distribution custom sequences rather than on the
independent LARD set. Finally, the style-transfer stage and the tracker-augmented
self-training variant are still ongoing, so the cross-technique comparison in
\Cref{tab:disc-summary} is not yet complete.
